\section{Related work}
\subsection{Tracking by detection and aerial MOT}
SORT established a compact online pipeline based on motion prediction and
assignment, while DeepSORT added appearance information to reduce identity
fragmentation \cite{bewley2016sort,wojke2017deepsort}. ByteTrack changed the
detector--tracker interface by associating high-confidence detections first and
then using lower-confidence boxes to recover objects that would otherwise be
lost \cite{zhang2022bytetrack}. Other modern trackers further refine motion and
appearance association. These trackers primarily improve the interpretation
of the boxes they receive; AC--MOT asks the complementary upstream question of
which observations, at which resolution and suppression policy, should reach
the tracker.

The broader tracking-by-detection literature includes detector-regression
tracking (Tracktor) \cite{bergmann2019tracktor}, joint detection and association
with transformer queries (TrackFormer) \cite{meinhardt2022trackformer}, and
end-to-end alternatives that change the association interface itself. These
lines of work motivate treating the detector--tracker boundary as an explicit
experimental variable rather than assuming that a tracker can compensate for
every detection-side choice. Standardized MOT evaluation and reporting are
rooted in MOTChallenge protocols \cite{milan2020motchallenge},
whereas the present aerial studies retain their project-specific protocol
boundaries.

VisDrone and UAVDT capture the small targets, changing target count, camera
motion, viewpoint variation, and clutter that make aerial MOT difficult
\cite{zhu2022visdrone,du2018uavdt}. In this paper they have separated roles:
historical VisDrone and UAVDT evidence describes earlier profiles, while a
modern frozen protocol tests our model. OATrack is a second tracker host rather
than a claim of superiority over its published system \cite{qi2026oatrack}; its
frozen repository implementation and tracker-side \texttt{min\_conf=0.40} are
kept distinct from the model's detector-side confidence floor.

\subsection{Detector operating points and adaptive inference}
Confidence floors, NMS overlap thresholds, and input resolution are often
treated as implementation settings. In aerial scenes they are coupled
scientific variables: resolution changes small-object evidence and cost,
confidence changes recall and clutter, and NMS changes whether nearby targets
remain separate. AC--MOT does not replace detector architecture research; it
makes these deployment choices explicit and lets them respond to scene evidence
without changing detector weights.

Dynamic inference systems such as Chameleon show that conditional computation
can adapt inference to content and resource constraints
\cite{jiang2018chameleon}. Resolution-adaptive networks and spatially adaptive
video models likewise use content to trade computation against quality
\cite{yang2020ranet}. AC--MOT is intentionally lighter: its
state uses scalar image statistics and prior tracker boxes, its action set is
finite, and its state machine is auditable line by line. It therefore tests an
operating-point intervention at the detector--tracker boundary rather than a
learned scheduler.

The detector host is also part of the evidence boundary. YOLOX provides a
representative modern anchor-free detector used by the U2MOT cross-pipeline
study \cite{ge2021yolox}, while U2MOT and SparseTrack illustrate that changes in
the association backend can alter the response to the same detector-side
policy \cite{liu2023u2mot,liu2025sparsetrack}. OATrack is therefore included as
a relevant recent UAV tracker host, not as a published system that this paper
claims to outperform \cite{qi2026oatrack}.

\subsection{Optimization and evaluation gap}
Black-box configuration is commonly implemented with automated hyperparameter
optimization; Optuna's define-by-run framework and TPE sampling are directly
relevant to the V1/V2 development record \cite{akiba2019optuna}. V2's explicit
trade-off is reported as a feasible Pareto front, using the standard notion of
non-dominance associated with multi-objective optimization \cite{deb2002nsga2}.
This distinction matters because V1 and V2 were not two arbitrary manual
settings: they answer different constrained questions.

MOTA emphasizes detection and identity errors, IDF1 emphasizes identity
precision and recall, and HOTA balances detection and association
\cite{luiten2021hota}. Their different responses are useful here because
selectivity may reduce FP and IDS while increasing FN. Sequence-level paired
bootstrap resampling preserves within-sequence dependence while quantifying
system differences \cite{efron1993bootstrap}.

The gap is consequently specific: rigorous evaluation of a training-free
online controller that changes detector resolution, confidence, and NMS in
aerial MOT while separating historical development from a frozen modern
protocol, testing two tracker hosts, measuring runtime, transferring without
retuning, and checking whether exact scene timing explains the gain. AC--MOT
addresses this gap empirically and does not claim to invent adaptive inference
or a new tracker.
